% ==================================================
% ABSTRACT
% ==================================================

\chapter*{Abstract}
\addcontentsline{toc}{chapter}{Abstract}

People read a great deal on the web but find it hard to return to a page
later, because browser history and bookmarks can only be searched by title
and address, and those are exactly the details people forget.
\textbf{MemoryLane} is a personal memory for the web, for students,
developers and researchers: pages are saved from the browser, with one click
or automatically, and found again later by describing what they were about.

A Chrome extension (Manifest V3) captures pages and YouTube captions; a
FastAPI backend cleans each page, splits it into 150-word overlapping
passages and embeds them with \texttt{bge-small-en-v1.5}; and Apache Solr
stores the passages for both keyword (BM25) and vector (KNN) retrieval. Every
search runs the two retrievers together and combines them with Reciprocal
Rank Fusion, understands site, time and address hints in the query, drops
weak matches with a relevance floor, and returns pages with readable snippets
and optional Gemini summaries. The system is multi-user, with Amazon Cognito
sign-in (PKCE) and per-user isolation enforced in the data-access layer, and
runs on one AWS EC2 instance with Docker Compose and Caddy at
\url{https://memorylane.site}.

The requirements are specified with unique identifiers and traced through
use cases, user stories, UML design models and twenty test cases, all of
which pass, together with 566 automated tests. On the SciFact benchmark the
live search puts the right abstract first for 71\% of claims and in the top
ten for 93\% (MRR@10 0.786, ahead of the prototype's 0.754), in about 76\,ms
per search; the evaluation also exposed a ranking defect whose fix raised the
first-result rate from 0.44 to 0.71. The main limitations are a scientific
benchmark that only partly reflects everyday browsing, a single server, and
page text that must be stored on that server.

\vspace{1em}
\noindent\textbf{Keywords:} personal information management, hybrid
retrieval, BM25, dense retrieval, Reciprocal Rank Fusion, browser extension,
software engineering
